\documentclass[11pt]{article}
\usepackage[margin=1in]{geometry}
\usepackage[pdftex]{graphicx}
\usepackage{amsmath,amssymb,amsthm,pgfplots,float}
\usepackage{minted}
\usepackage{custom}

% color boxes
\usepackage[many]{tcolorbox}
\tcbuselibrary{theorems}
\newtcbtheorem{psexample}{Example}{colback=green!10,colframe=green!40!black,fonttitle=\bfseries,breakable,enhanced}{ex}
\newtcbtheorem{psproblem}{Problem}{colback=blue!10,colframe=blue!40!black,fonttitle=\bfseries,breakable,enhanced}{id}
\newtcbtheorem{psremark}{Remark}{colback=purple!10,colframe=purple!40!black,fonttitle=\bfseries,breakable,enhanced}{rm}
\newtcbtheorem{pssolution}{Solution}{colback=orange!10,colframe=orange!40!black,fonttitle=\bfseries,breakable,enhanced}{sn}
\tcbset{noparskip/.style={before={\par\pagebreak[0]\medskip\parindent=0pt},
after={\par\medskip}}}

% formatting for problems and solutions
\definecolor{ptred}{rgb}{0.7,0.1,0.1}
\newcommand{\ptfmt}[1]{\textbf{\color{ptred}#1\color{black}}}
\definecolor{advblue}{rgb}{0.1,0.1,0.7}
\newcommand{\advanced}{\!\textbf{\color{advblue}[A]\color{black}}\ }
\newcommand{\blue}[1]{\textcolor{blue}{#1}}
\newcommand{\orange}[1]{\textcolor{orange}{#1}}

\newtheorem*{solution}{Solution}
\newtheorem*{answer}{Answer}

\makeatletter
\newtheoremstyle{mystyle}
{\topsep}               % space above
{\topsep}               % space below
{}                      % bodyfont
{}                      % indent
{\bfseries}             % headfont
{}                      % punctuation
{0.6em}                 % space after head
{\llap{\hspace{.6em}}\thmname{#1}\thmnumber{ #2}\thmnote{\normalfont{ (#3)}}{\bfseries .}}  %theoremheadspec
\theoremstyle{mystyle}
\newtheorem{pproblem}{Problem}
\makeatother

% clocks for time limits
\usepackage{clock} 
% headers and footers
\usepackage{fancyhdr}
\usepackage{xcolor}
\pagestyle{fancy}
\lhead{Filip Ramazan}
\chead{}
\rhead{STAT 760}
\lfoot{}
\cfoot{\thepage}
\rfoot{}
\renewcommand{\headrulewidth}{0.4pt}
\setlength{\headheight}{14pt}

\newcommand{\psettitle}[1]{
    \begin{center}
    \huge \textbf{#1}
    \end{center}
}

\linespread{1.03} % give a little extra room
\setlength{\parindent}{0.2in} % reduce paragraph indent a bit

\usemintedstyle{paraiso-dark}
\definecolor{bg}{rgb}{0.13,0.13,0.13}
\definecolor{fg}{rgb}{0.95,0.95,0.95}
\setmintedinline{bgcolor=bg}
\DeclareUnicodeCharacter{03BB}{\ensuremath{\lambda}}

\begin{document}
\psettitle{Midterm}
\begin{center}
{\large March $20^{\text{th}}$, 2026}
\end{center}
I began my solution by implementing the algorithm given in the instructions. I chose hyperparameters after running a program to test different ones, but I'll detail that at the end. I use a RBF (Gaussian) kernel.
At first, I chose a multiplicative update. The main issues I found with multiplicative updates were that: 1. multiplicative updates would stall convergence, and most alphas would get near zero, a multiplicative update cannot revive it easily. 
2. they would need scaling to maintain the constraints, and that would typically mean needing many, many iterations to converge. Also, I tried, but failed, to add things like a learning rate schedule, as the updates became too small too quick.  For additive updates, it was easier to handle with the constraints in this case. I initially tried to compute the violation $\sum{\alpha_i \cdot y_i}$, and if the violation $> 0$ I reduced all positive-class alphas uniformly by the sum divided by the number of points that wouldn't become negative after updating (and vice versa for negative violation).
I was getting decent results with this method, but it required 6,000,000 iterations. For measuring error, I created a test set of 2000 random data points, and also estimated the area between the known boundary and my model by using the sum of the areas of hundreds of trapezoids. The area estimate helps me quantify the shape/form of the model, and the test set the effectiveness.
\begin{minted}[
    bgcolor=bg,
    fontfamily=tt,
    fontsize=\small,
    linenos,
    breaklines
]{python}             
Dual objective               62.9115
Training accuracy             1.0000
Support vectors                   25
Max alpha                    13.4988
Boundary area error           0.0241
Test accuracy (n=2000)        0.9745
\end{minted}
My method up until now is clearly inefficient. The learning rate is constant (at $\gamma=0.002$), and small. For meaningful changes, I need a learning rate that updates more for the larger violators and smaller for the smaller ones. I would also benefit from a method that updates only the worst alphas, rather than choosing randomly like I was doing up until now, and possibly not a constraint repair system that can slightly move me off the right objective ascent path. To set this up, I shifted to a soft margin SVM too. I noticed that in the results for the first version, some of the randomly generated data was clumped, and the boundary tended toward it there. See here:
\begin{figure}[H]
    \centering
    \includegraphics[width=0.5\textwidth]{/Users/filipr/Library/CloudStorage/OneDrive-DavidsonAcademy/STAT 760 - Rojas/svm_stage1.png}
\end{figure}
\noindent Soft margin, capping the influence of any single one point, will hopefully generalize better. For this, I derived a formula for the learning rate as a function of the gradients, using the Lagrangian:
\[
L(\alpha) = \sum_i \alpha_i - \frac{1}{2} \sum_{i,j} \alpha_i \alpha_j y_i y_j K_{ij}, \quad \sum_i \alpha_i y_i = 0
\]

\[
\text{Choose a pair } k,l \text{ with } y_k = 1, y_l = -1, \quad \alpha(\delta) = \alpha + \delta (e_k + e_l)
\]

\[
L(\delta) = \sum_i \alpha_i + 2\delta - \frac{1}{2} (\alpha + \delta(e_k+e_l))^\top Q (\alpha + \delta(e_k+e_l))
\]

\[
L(\delta) = \sum_i \alpha_i - \frac{1}{2} \alpha^\top Q \alpha + 2\delta - \delta (Q\alpha)\cdot k - \delta (Q\alpha)\cdot l - \frac{1}{2} \delta^2 (Q\cdot {kk} + Q\cdot {ll} - 2Q_{kl})
\]

\[
\frac{d L}{d\delta} = 2 - (Q\alpha)\cdot k - (Q\alpha)\cdot l - \delta (Q\cdot {kk}+Q\cdot {ll}+2Q_{kl}) = g_k + g_l - \delta (Q_{kk}+Q_{ll}+2Q_{kl})
\]

\[
0 = g_k + g_l - \delta (Q_{kk}+Q_{ll}+2Q_{kl}) \quad \Rightarrow \quad \delta^* = \frac{g_k + g_l}{Q_{kk}+Q_{ll}+2Q_{kl}}
\]
\noindent Now, translate this in terms of the kernel.
\[
Q_{kk} + Q_{ll} + 2 Q_{kl} = K_{kk} + K_{ll} - 2 K_{kl} \quad \Rightarrow \quad \delta^* = \frac{g_k + g_l}{K_{kk} + K_{ll} - 2 K_{kl}}
\]
I used this formula which helps me maximize the change to the Lagrangian with every update to each pair. For picking the pairs, I focus only on the most violating pair for one-half of the updates (odd iterations), and the other half are are chosen at random (even iterations). This way I ensure the $\alpha$s are maximized for the points close to the boundary and the others go to zero. For the soft margin, the clipping to maintain alpha under C is applied when calculating delta. Here are the results from this method:
 \begin{minted}[
    bgcolor=bg,
    fontfamily=tt,
    fontsize=\small,
    linenos,
    breaklines
]{python}                     
Dual objective               37.9882
Training accuracy             1.0000
Support vectors                   31
Max alpha                     4.0000
Bias                          0.0123
Boundary area error           0.0200
Test accuracy (n=2000)        0.9805
\end{minted}
These results are clearly better, with both boundary and test accuracy improving. There was a modest increase in support vectors. The main problem I see with this implementation, though, is that the Lagrangian improves only by hundredths after the first 10,000 iterations. To fix this, I think I should focus on choosing the pair with the maximum gain in Lagrangian, rather than focusing on the pair that each have the largest absolute gradients. To do this, I derived a formula:
\[
L(\alpha)=\sum_i \alpha_i-\frac12\sum_{i,j}\alpha_i\alpha_j y_i y_j K_{ij}
\]

\[
\alpha_k'=\alpha_k+\delta,\qquad \alpha_l'=\alpha_l+\delta
\]

\[
L(\delta)=L(\alpha)+\delta(g_k+g_l)-\frac12\delta^2(K_{kk}+K_{ll}-2K_{kl})
\]

\[
g_i=1-y_i\sum_j \alpha_j y_j K_{ji}
\]

\[
\frac{dL}{d\delta}=g_k+g_l-(K_{kk}+K_{ll}-2K_{kl})\delta
\]

\[
0=g_k+g_l-(K_{kk}+K_{ll}-2K_{kl})\delta^*
\]

\[
\delta^*=\frac{g_k+g_l}{K_{kk}+K_{ll}-2K_{kl}}
\]
\[
\begin{aligned}
\Delta L &= L(\delta^*)-L(0)=\delta^*(g_k+g_l)-\frac12(K_{kk}+K_{ll}-2K_{kl})(\delta^*)^2 \\
&= \frac{g_k+g_l}{K_{kk}+K_{ll}-2K_{kl}}(g_k+g_l)
 - \frac12 (K_{kk}+K_{ll}-2K_{kl})
 \left(\frac{g_k+g_l}{K_{kk}+K_{ll}-2K_{kl}}\right)^2 \\
&= \frac{(g_k+g_l)^2}{K_{kk}+K_{ll}-2K_{kl}}
 - \frac12 \frac{(g_k+g_l)^2}{K_{kk}+K_{ll}-2K_{kl}} = \frac{(g_k+g_l)^2}{2\bigl(K_{kk}+K_{ll}-2K_{kl}\bigr)}
\end{aligned}
\]
\noindent Now, I used this formula in my code, I extracted the top 15 points, with the highest gradients for positive and negative. I compute the expected gain for each, and then choose the best pair. The even iterations remained random points. My final results were marginally better on the test set:
 \begin{minted}[
    bgcolor=bg,
    fontfamily=tt,
    fontsize=\small,
    linenos,
    breaklines
]{python}
Dual objective               33.9284
Training accuracy             1.0000
Support vectors                   33
Max alpha                     3.0000
Bias                          0.0107
Boundary area error           0.0202
Test accuracy (n=2000)        0.9825
\end{minted}
\begin{figure}[H]
    \centering
    \includegraphics[width=0.5\textwidth]{/Users/filipr/Library/CloudStorage/OneDrive-DavidsonAcademy/STAT 760 - Rojas/svm_stage3.png}
\end{figure}
As I didn't have a larger improvement, I stopped trying new methods. Code is shown below for the last method I tried, and it has in it a loop that tests various combinations of hyperparameters. As for possible other improvements or directions I could taken, I could've:
\begin{itemize}
\item Temporarily removed points that couldn't become support vectors, like those whose $\alpha_i$ have been pinned at 0 or C for many consecutive iterations, which are unlikely to change. 
\item I could make it quite faster by using broadcasting for my kernel matrix computation.
\item The working set size $Z=15$ was fixed. An adaptive Z that grows when gains are small and shrinks when they're large could balance exploration and exploitation better, especially in the later iterations where the Lagrangian improves only by hundredths.
\item  The values of $\gamma$, $C$, and $Z$ were chosen by running a grid search on test accuracy alone. $k$-fold cross-validation on the training set would give a less biased estimate of generalization, and might have suggested different hyperparameters.
\end{itemize}
 \begin{minted}[
    bgcolor=bg,
    fontfamily=tt,
    fontsize=\small,
    linenos,
    breaklines
]{python}
import numpy as np
from svm_utils import (
    generate_data, kernel_matrix, init_alphas,
    dual_objective, compute_bias, training_accuracy,
    test_accuracy, boundary_area_error, plot_results,
)
gamma    = 0.8 # kernel width
n_points = 150 # training set size
C_svm    = 3.0 # soft margin upper bound on alphas
max_iter = 1_000_000 # number of pairwise iterations
Z        = 15 # working-set size per class

def optimize_v3(x, y, K, Q, max_iter, Z=5, C_svm=10.0):
    alpha = init_alphas(y)
    history = []
    pos_idx = np.where(y == 1)[0] # indices of positive-class points
    neg_idx = np.where(y == -1)[0] # indices of negative-class points

    for t in range(1, max_iter + 1):
        # compute all classifications and gradients at once
        C_all = (alpha * y) @ K
        grad_all = 1.0 - y * C_all

        if t % 2 == 1:
            # odd iterations: pick the pair with largest gain in Lagrangian from working set
            pos_grads = grad_all[pos_idx]
            neg_grads = grad_all[neg_idx]
            # get working set of Z points with largest gradients
            top_pos = pos_idx[np.argsort(pos_grads)[-Z:]]
            top_neg = neg_idx[np.argsort(neg_grads)[-Z:]]
            best_gain = -1.0
            best_pair = (top_pos[-1], top_neg[-1])
            # iterate through all pairs in working set
            for i in top_pos:
                for j in top_neg:
                    denominator = K[i, i] + K[j, j] - 2.0 * K[i, j]
                    if denominator > 1e-12:
                        predicted_gain = (grad_all[i] + grad_all[j])**2 / (2.0 * denominator)
                        if predicted_gain > best_gain:
                            best_gain = predicted_gain
                            best_pair = (i, j)
            k, l = best_pair
        else:
            # EVEN: pick a random pair (ensures all points get visited)
            k = pos_idx[np.random.randint(len(pos_idx))]
            l = neg_idx[np.random.randint(len(neg_idx))]

        # optimal delta:
        # by plugging alpha_k += delta, alpha_l += delta into the
        # dual Lagrangian and setting dL/d(delta) = 0, we get:
        # delta = (grad_k + grad_l) / (K_kk + K_ll - 2*K_kl)
        # this takes the analytically optimal step for this pair
        optimal_lr = K[k, k] + K[l, l] - 2.0 * K[k, l]
        if optimal_lr < 1e-12:
            continue
        delta = (grad_all[k] + grad_all[l]) / optimal_lr
        # clip delta so both alphas stay in [0, C_svm]
        if delta > 0:
            delta = min(delta, C_svm - alpha[k], C_svm - alpha[l])
        else:
            delta = max(delta, -alpha[k], -alpha[l])
        if abs(delta) < 1e-12:
            continue
        # apply the update
        alpha[k] += delta
        alpha[l] += delta
        # zero out negligible alphas
        if alpha[k] < 1e-8:
            alpha[k] = 0.0
        if alpha[l] < 1e-8:
            alpha[l] = 0.0
        # log periodically
        if t % 1000 == 0:
            b = compute_bias(alpha, y, K, C_svm)
            L = dual_objective(alpha, Q)
            acc = training_accuracy(alpha, y, K, b)
            n_sv = np.sum(alpha > 1e-4)
            history.append((t, L, acc, n_sv))
            print(f"  iter {t:>7d}   L = {L:>10.4f}   acc = {acc:.4f}   SVs = {n_sv}")
    return alpha, history

if __name__ == '__main__':
    np.random.seed(16)
    x, y, _ = generate_data(n_points, seed=16)
    n = len(y)
    n_pos = np.sum(y == 1)
    n_neg = np.sum(y == -1)
    print(f"Training points:  {n}  (+1: {n_pos}, -1: {n_neg})")

    # define hyperparameter grid
    param_grid = {
        'gamma': [0.1, 0.5, 0.8, 1.0],
        'Z': [5, 10, 15],
        'C_svm': [1.0, 3.0, 5.0, 10.0]
    }

    # run grid search
    best_params, best_alpha, best_bias, best_acc = grid_search(
        x, y, optimize_v3, param_grid, max_iter=max_iter
    )

    print("Plotting best model...")
    best_gamma = best_params.get('gamma', 0.8)
    plot_results(x, y, best_alpha, best_gamma, bias=best_bias, title=f"Best SVM: {best_params}", filename='svm_v3_best.png')

    bias = compute_bias(best_alpha, y, kernel_matrix(x, best_gamma), best_params.get('C_svm', 10.0))
    acc = training_accuracy(best_alpha, y, kernel_matrix(x, best_gamma), bias=bias)
    Q = (y[:, None] * y[None, :]) * kernel_matrix(x, best_gamma)
    lagrangian = dual_objective(best_alpha, Q)
    sv = np.sum(best_alpha > 1e-4)

    print(f"\n{'Best Model Results':^60}")
    print(f"{'-' * 60}")
    print(f"{'Parameters':<25} {str(best_params)}")
    print(f"{'Dual objective':<25} {lagrangian:>10.4f}")
    print(f"{'Training accuracy':<25} {acc:>10.4f}")
    print(f"{'Support vectors':<25} {sv:>10d}")
    print(f"{'Max alpha':<25} {np.max(best_alpha):>10.4f}")
    print(f"{'Bias':<25} {bias:>10.4f}")
    area = boundary_area_error(x, y, best_alpha, best_gamma, bias=bias)
    print(f"{'Boundary area error':<25} {area/25:>10.4f}")
    tacc = test_accuracy(x, y, best_alpha, best_gamma, bias=bias)
    print(f"{'Test accuracy (n=2000)':<25} {tacc:>10.4f}")
    print(f"{'-' * 60}")
 \end{minted}
For each of my versions, I used these "util" functions:
 \begin{minted}[
    bgcolor=bg,
    fontfamily=tt,
    fontsize=\small,
    linenos,
    breaklines
]{python}
import numpy as np
import matplotlib.pyplot as plt

def rbf_kernel(xi, xj, gamma):
    # K(xi, xj) = exp(-gamma * ||xi - xj||^2)
    diff = xi - xj
    return np.exp(-gamma * np.dot(diff, diff))

def kernel_matrix(x, gamma):
    # precompute full n x n kernel table
    n = x.shape[0]
    K = np.zeros((n, n))
    for i in range(n):
        for j in range(i, n):
            K[i, j] = rbf_kernel(x[i], x[j], gamma)
            K[j, i] = K[i, j]
    return K

def generate_data(n_points, seed=16):
    # +1 lies below sin(x) + 2.5
    # -1 lies above sin(x) + 2.5
    np.random.seed(seed)
    x = np.random.uniform(0, 5, size=(n_points, 2))
    boundary = np.sin(2 * np.pi * x[:, 0] / 5.0) + 2.5
    y = np.where(x[:, 1] < boundary, 1.0, -1.0)
    return x, y, boundary

def classify(alpha, y, K_col):
    # C(x) = sum_i alpha_i * y_i * K(x_i, x)
    # K_col is the column of the kernel matrix for the query point
    return np.sum(alpha * y * K_col)

def init_alphas(y):
    # start with small random positive alphas
    n = len(y)
    alpha = np.random.uniform(0.001, 0.01, size=n)
    # enforce sum(alpha_i * y_i) = 0 by adjusting each class
    # so that total positive-class alpha equals total negative-class alpha
    pos = (y == 1)
    neg = (y == -1)
    sum_pos = np.sum(alpha[pos])
    sum_neg = np.sum(alpha[neg])
    # scale the larger side down to match the smaller
    if sum_pos > sum_neg:
        alpha[pos] *= sum_neg / sum_pos
    else:
        alpha[neg] *= sum_pos / sum_neg
    return alpha


def repair_constraint(alpha, y, tol=1e-12):
    violation = np.sum(alpha * y)
    if abs(violation) < tol:
        return alpha
    # determine which class to reduce
    if violation > 0:
        mask = (y == 1)
    else:
        mask = (y == -1)
    remaining = mask.copy()
    while True:
        idx = np.where(remaining)[0]
        if len(idx) == 0:
            break
        adjustment = abs(violation) / len(idx)
        alpha[idx] -= adjustment
        # detect which ones went negative
        neg = alpha[idx] < 0
        if not np.any(neg):
            break
        # clip them
        alpha[idx[neg]] = 0.0
        # remove them from redistribution pool
        remaining[idx[neg]] = False
        # recompute violation after clipping
        violation = np.sum(alpha * y)
        if abs(violation) < tol:
            break
    return alpha

def dual_objective(alpha, Q):
    # L(alpha) = sum(alpha) - 0.5 * alpha^T Q alpha
    return np.sum(alpha) - 0.5 * alpha @ Q @ alpha

def compute_bias(alpha, y, K, C_svm):
    # bias is computed from "free" support vectors: 0 < alpha_k < C
    # for these points, y_k * (C(x_k) + b) = 1
    # so b = y_k - C(x_k)
    tol = 1e-4
    free_sv = np.where((alpha > tol) & (alpha < C_svm - tol))[0]
    if len(free_sv) == 0:
        # fallback: use all support vectors
        free_sv = np.where(alpha > tol)[0]
    if len(free_sv) == 0:
        return 0.0
    # b = average of (y_k - C(x_k)) over free SVs
    b_values = []
    for k in free_sv:
        C_k = classify(alpha, y, K[k, :])
        b_values.append(y[k] - C_k)
    return np.mean(b_values)

def training_accuracy(alpha, y, K, bias=0.0):
    n = len(y)
    correct = 0
    for i in range(n):
        C_i = classify(alpha, y, K[i, :]) + bias
        pred = np.sign(C_i)
        if pred == 0:
            pred = 1.0
        if pred == y[i]:
            correct += 1
    return correct / n

def test_accuracy(x_train, y_train, alpha, gamma, bias=0.0,
                  n_test=2000, test_seed=999):
    # generate a fresh test set and evaluate the model on it
    x_test, y_test, _ = generate_data(n_test, seed=test_seed)
    # compute kernel between training points and test points
    # K_test[i, j] = K(x_train[i], x_test[j])
    diff = x_train[:, np.newaxis, :] - x_test[np.newaxis, :, :]
    sq_dist = np.sum(diff ** 2, axis=2)
    K_test = np.exp(-gamma * sq_dist)
    # f(x_test_j) = sum_i alpha_i y_i K(x_train_i, x_test_j) + bias
    f_test = (alpha * y_train) @ K_test + bias
    preds = np.sign(f_test)
    preds[preds == 0] = 1.0
    return np.mean(preds == y_test)

def boundary_area_error(x, y, alpha, gamma, bias=0.0, n_grid=1000):
    # compute area between SVM decision boundary and true sine curve
    # by finding f(x1, x2) = 0 for each x1 and comparing to sin boundary
    x1_vals = np.linspace(0, 5, n_grid)
    x2_fine = np.linspace(0, 5, n_grid)
    svm_boundary = np.full(n_grid, np.nan)
    for i, x1 in enumerate(x1_vals):
        # evaluate f at all x2 values for this x1
        pts = np.column_stack([np.full(n_grid, x1), x2_fine])
        diff = x[:, np.newaxis, :] - pts[np.newaxis, :, :]    # (n, n_grid, 2)
        sq_dist = np.sum(diff ** 2, axis=2)
        K_pts = np.exp(-gamma * sq_dist)
        f_vals = (alpha * y) @ K_pts + bias
        # find zero crossing (sign change)
        for j in range(n_grid - 1):
            if f_vals[j] * f_vals[j + 1] <= 0:
                # linear interpolation to find the zero
                t = f_vals[j] / (f_vals[j] - f_vals[j + 1])
                svm_boundary[i] = x2_fine[j] + t * (x2_fine[j + 1] - x2_fine[j])
                break
    # true boundary
    true_boundary = np.sin(2 * np.pi * x1_vals / 5.0) + 2.5
    # only use points where we found the SVM boundary
    valid = ~np.isnan(svm_boundary)
    if np.sum(valid) < 2:
        return float('inf')
    # area = approximate integral of |svm - true| over [0, 5]
    area = np.trapezoid(np.abs(svm_boundary[valid] - true_boundary[valid]),
                        x1_vals[valid])
    return area

def plot_results(x, y, alpha, gamma, bias=0.0,
                 title='Kernel SVM',
                 filename='svm_result.png'):
    _, ax = plt.subplots(figsize=(8, 6))
    # scatter training points
    pos = (y == 1)
    neg = (y == -1)
    ax.scatter(x[pos, 0], x[pos, 1], c='blue', marker='o', s=20, label='+1')
    ax.scatter(x[neg, 0], x[neg, 1], c='red',  marker='x', s=20, label='-1')
    # circle support vectors (alpha > 0.5)
    sv = np.where(alpha > 0.5)[0]
    if len(sv) > 0:
        ax.scatter(x[sv, 0], x[sv, 1], facecolors='none', edgecolors='black',
                   s=120, linewidths=1.5, label=f'SV (alpha > 0.5, n={len(sv)})')
    # decision boundary, margins on a grid
    xx, yy = np.meshgrid(np.linspace(0, 5, 300), np.linspace(0, 5, 300))
    grid = np.c_[xx.ravel(), yy.ravel()]
    # compute f(x) = C(x) + bias for each grid point
    diff = x[:, np.newaxis, :] - grid[np.newaxis, :, :]
    sq_dist = np.sum(diff ** 2, axis=2)
    K_grid = np.exp(-gamma * sq_dist)
    f_grid = (alpha * y) @ K_grid + bias
    f_grid = f_grid.reshape(xx.shape)
    # decision boundary (f = 0) and margins (f = +-1)
    ax.contour(xx, yy, f_grid, levels=[-1], colors='gray', linestyles='--', linewidths=0.8)
    ax.contour(xx, yy, f_grid, levels=[0],  colors='black', linewidths=1.5)
    ax.contour(xx, yy, f_grid, levels=[1],  colors='gray', linestyles='--', linewidths=0.8)
    # true sine boundary
    xs = np.linspace(0, 5, 300)
    ax.plot(xs, np.sin(2 * np.pi * xs / 5.0) + 2.5, 'g-', linewidth=1.5, label='true boundary')
    ax.set_xlim(0, 5)
    ax.set_ylim(0, 5)
    ax.set_xlabel('x_1')
    ax.set_ylabel('x_2')
    ax.set_title(title)
    ax.legend(loc='upper right', fontsize=8)
    plt.tight_layout()
    plt.savefig(filename, dpi=150)
    plt.show()
    print(f"Plot saved to {filename}")

def grid_search(x, y, optimize_fn, param_grid, max_iter):
    keys = list(param_grid.keys())
    values = param_grid.values()
    combinations = list(itertools.product(*values))
    best_params = None
    best_alpha = None
    best_acc = -1.0
    best_bias = 0.0
    print(f"{'Grid Search':^60}")
    print("="*60)
    
    for combo in combinations:
        params = dict(zip(keys, combo))
        gamma = params.get('gamma', 0.8)
        Z = params.get('Z', 5)
        C_svm = params.get('C_svm', 10.0)
        
        K = kernel_matrix(x, gamma)
        Q = (y[:, None] * y[None, :]) * K
        
        print(f"Testing parameters: {params}...")
        alpha, history = optimize_fn(x, y, K, Q, max_iter=max_iter, Z=Z, C_svm=C_svm)
        bias = compute_bias(alpha, y, K, C_svm)
        acc = test_accuracy(x, y, alpha, gamma, bias=bias)
        print(f"Result for {params}: Test Acc = {acc:.4f}")
        print("-" * 60)
        if acc > best_acc:
            best_acc = acc
            best_params = params
            best_alpha = alpha
            best_bias = bias
    print("="*60)
    print(f"Best parameters: {best_params} with test acc: {best_acc:.4f}")
    
    return best_params, best_alpha, best_bias, best_acc
\end{minted}
\end{document}